
\documentclass[%
 reprint,
superscriptaddress,
%groupedaddress,
%unsortedaddress,
%runinaddress,
%frontmatterverbose, 
%preprint,
%preprintnumbers,
%nofootinbib,
%nobibnotes,
%bibnotes,
aps,
pra,
onecolumn, 
%prb,
%rmp,
%prstab,
%prstper,
%floatfix,
]{revtex4-1}

\usepackage{dcolumn}% Align table columns on decimal point
\usepackage{amsmath,amssymb,graphicx,comment,bm,color}
%\usepackage[pagebackref]{hyperref}
\usepackage{tikz}

\begin{document}

\section{Summary}


The numerical code solves two types of one-dimensional optical wave turbulence equation. The numerical code solves for the 1D complex wave function $\psi = \psi(x,t)$ with periodic boundaries of size $L$.

The first type we call the non-local optical wave turbulence equation

\begin{align*}
i\frac{\partial \psi}{\partial t}=-\frac{1}{2}\frac{\partial^2 \psi}{\partial x^2} -\frac{1}{2}\psi\left( 1- \frac{1}{g}\frac{\partial^2}{\partial x ^2}\right)^{-1}|\psi|^2 -iD + iF
\end{align*}
where $F$ is an additive forcing and the penultimate last two terms are optional dissipation terms.


The system is Hamiltonian in the no forcing and no dissipation limit ($F=D=0s$). It conserve the energy or Hamiltonian given by:

\begin{align*}
H = H_{lin}+H_{non} =   \int \frac{1}{2}\left| \frac{\partial \psi}{\partial x}\right|^2   - \frac{1}{4}\left[ \left( 1- \frac{1}{g}\frac{\partial^2}{\partial x^2}\right)^{-1/2}|\psi|^2\right]^2 \ dx.
\end{align*}

The second type we call the long-wave optical wave turbulence equation which is derived as the leading order expression of the non-local equation in the limit $(1/g)\partial^2/\partial x^2 \ll 1$:

\begin{align*}
i\frac{\partial \psi}{\partial t}=-\frac{1}{2}\frac{\partial^2 \psi}{\partial x^2}  -\frac{1}{2}\psi|\psi|^2 + \frac{1}{2g}\psi\frac{\partial^2|\psi|^2}{\partial x ^2} -iD + iF
\end{align*}
where $F$ is an additive forcing and the penultimate last two terms are optional dissipation terms.


The system is Hamiltonian in the no forcing and no dissipation limit ($F=D=0s$). It conserve the energy or Hamiltonian given by:

\begin{align*}
H = H_{lin}+H_{non1} + H_{non2} =   \int \frac{1}{2}\left| \frac{\partial \psi}{\partial x}\right|^2   - \frac{1}{4}|\psi|^4 + \frac{1}{4g} \left(\frac{\partial |\psi|^2}{\partial x}\right)^{2} \ dx.
\end{align*}

\subsection{Parameters}

The main dynamical parameters are the chemical potential $\mu$ and level of nonlinearity $g$.

The structure of the additive forcing is such that
\begin{align*}
F(x,t) = \sum_{k} f_k \eta_k(t),
\end{align*}
where $\eta_k(t)$ is a complex Gaussian white noise at mode $k$ and $f_k\in \mathbb{R}$ is the distribution of the forcing amplitudes in Fourier space.

The dissipation is of the form of hypo- an hyper-viscosity terms
\begin{align*}
D(x,t) = \alpha \left(-\frac{\partial^2}{\partial x^2}\right)^{\alpha_{power}}\psi +\nu \left(-\frac{\partial^2}{\partial x^2}\right)^{\nu_{power}}\psi
\end{align*}
where the parameters $\alpha$ and $\nu$ correspond to the hypo- an hyper-viscosity coefficients with the order of the two dissipation terms given by $\alpha_{power}\leq 0$ and $\nu_{power}>0$.

\subsection{Discretization}

The equation is solved by a pseudo-spectral spatial method in a periodic 1D domain, with full dealiasing with the 3/2-rule. We use the FFTW library that stores the wavenumbers in the for $\mathbf{k} = 0,1,2,\dots, N/2, -N/2 +1, \dots, -2,-1$, where $N$ is the spatial resolution.

Time-stepped using an fourth-order exponential time-differencing Runge-Kutta method (see~\cite{Kassam_fourth-order_2005,cox_exponential_2002} for details).


\section{Installation}

The code uses the following libraries that will need to be installed by the user
\begin{itemize}
\item gcc (compiler)
\item fftw-3 (fast Fourier transform library)
\item armadillo (C++ linear algebra library)
\item openBLAS (linear algebra library that is used by armadillo)
\end{itemize}

All of the above can be installed via the local repositories on Linux, or via Macports (and possibly Homebrew) on macOS.


\section{How to Run the Code}

The numerical code contains a \texttt{/src/Makefile} that will compile the code using the gcc compiler and create an executable \texttt{/src/owt1d}. At the very top are paths to the default locations of the linked libraries for both Linux and MacOS. Simply comment out the one not in use. Bear in mind that the location of the libraries may differ on your local machine.

All global parameters are stored in the header file \texttt{/src/const.h}. Any changes in \texttt{/src/const.h} will require a recompilation of the numerical code before running.

Once \texttt{/src/owt1d} is created, it can be moved to any directory one wishes, but the \texttt{/data/} folder must be present in the running directory.

\subsection{The \texttt{/data/} Folder}
To run the code, the user \textbf{must} create a \texttt{/src/data/} directory in the same folder as the created executable.  In the \texttt{/data/} folder the user must create a file called \texttt{/data/curframe.dat} that will contain two numbers separated by a tab. The numbers represent the current time and the label of the last generated file number. Each time the code outputs a new file for the data, \texttt{/data/curframe.dat} will be automatically updated with the new time and label of the last data file. The code uses this information to when restarting from the last generated file.

\begin{center}
\begin{tabular}{c|c}
current time & file number \\
\hline
$z$ & \texttt{XXXXXX} \\

\end{tabular}
\end{center}

The folder \texttt{./data/} will also contain all the data for the wave function $\psi$.  They are stored in files called \texttt{/data/psi.XXXXXX} where \texttt{XXXXXX} is the file number padded by zeros. Files \texttt{/data/psi.XXXXXX} are core data for the numerical code.

\begin{center}
\begin{tabular}{c|c|c}
$z$ & $Re(\psi)$ & $Im(\psi)$\\
\hline
$\vdots$ & $\vdots$ & $\vdots$
\end{tabular}
\end{center}

\subsection{The \texttt{./output/} Folder}

The code will automatically generate (if not already present) another folder called \texttt{./output/} that will contain all necessary in-code post-processing of the data. Usually, this includes the data used for debugging and verification. The code will output files at the same intervals as the main data \texttt{/data/psi.XXXXXX}.

\texttt{/output/energy.XXXXXX} records the Hamiltonian or energy of the equation to ensure that it is properly conserved by the dynamics. It is recorded in the form of a row of four numbers:
\begin{center}
\begin{tabular}{c|c|c|c}
current time & linear energy   & nonlinear energy  & total (linear + nonlinear) energy\\
\hline
$z$ & $H_{lin}$& $H_{non}$ & $H_{lin}+H_{non}$
\end{tabular}
\end{center}

\texttt{/output/wave.XXXXXX} records the Hamiltonian or energy of the equation to ensure that it is properly conserved by the dynamics. It is recorded in the form of a row of four numbers:
\begin{center}
\begin{tabular}{c|c|c}
current time & total wave action & zeroth mode of wave action\\
\hline
z & $N$ & $|\hat{\psi}_{k=0}|^2$
\end{tabular}
\end{center}


\texttt{/output/spec.XXXXXX} records the Fourier wave action spectrum of $|\hat{\psi}_k|^2$ and wave energy spectrum  $\omega_k|\hat{\psi}_k|^2 = k^2|\hat{\psi}_k|^2$ vs $k$. It is recorded in the form of a five columns of numbers:
\begin{center}
\begin{tabular}{c|c|c|c|c}
$|k|$ & $|\hat{\psi}_k|^2$ & $k^2|\hat{\psi}_k|^2$ & $\langle |\hat{\psi}_k|^2 \rangle $ & $\langle k^2|\hat{\psi}_k|^2 \rangle$\\
\hline
$\vdots$ & $\vdots$ & $\vdots$ & $\vdots$ & $\vdots$
\end{tabular}
\end{center}




\texttt{/output/flux.XXXXXX} records spectral Fourier flux vs $k$ for both the wave action and energy (instantaneous as well as time averaged over data since start of code). It is recorded in the form of a five columns of numbers:
\begin{center}
\begin{tabular}{c|c|c|c| c}
$|k|$ & wave action flux $\eta_k$ & energy flux $\epsilon_k$ & $\langle \eta_k \rangle $ & $\langle \epsilon_k\rangle$\\
\hline
$\vdots$ & $\vdots$ & $\vdots$ & $\vdots$ & $\vdots$
\end{tabular}
\end{center}


\text{/output/dis.XXXXXX} records the total wave action and energy dissipated via the terms that constitute $D$ at the time of output. It is recorded in the form of a row of five  numbers:
\begin{center}
\begin{tabular}{c|c|c|c|c}
time  & wave action dissipated by $\alpha$& wave action dissipated by $\nu$ & energy dissipated by $\alpha$  & energy dissipated by $\nu$\\
\hline
. & .& .& .& 
\end{tabular}
\end{center}



Using terminal command \texttt{cat energy.* > energy.tot} will append all individual output files into one single file for plotting.


\subsection{The \texttt{./initial/} Folder}

As probably surmised, the code initializes from the data file \texttt{/data/psi.XXXXXX} stored in \texttt{/data/} to which the file number in \texttt{/data/curframe.dat} points. Therefore, to initialise the code when not restarting one needs to include an initial data file, normally \text{/src/data/psi.000000} (it requires for \texttt{/data/curframe.dat} to be set to $0\quad 0$). 

If you require a zero-state initial condition because you are running a forced/dissipated simulation, you can set \texttt{XXXXXX} in \texttt{/data/curframe.dat} to be any negative integer and the code will interpret this as you wanting to create a zero-state initial condition and will generate this automatically.

If you want to generate a custom (bespoke) initial condition, then next to the \texttt{/src/} folder is the \texttt{/initial/} folder that contains a simple .cpp file for generating a custom initial state that can be copied into \texttt{/data/}.










\bibliography{library}

\end{document}